\exercisesection{}

\begin{enumerate}
\def\labelenumi{\arabic{enumi})}
\tightlist
\item
  设 G 为紧拓扑群，B 为 G 的不可约酉表示的特征标所成的族。
\end{enumerate}

\begin{enumerate}
\def\labelenumi{\alph{enumi})}
\item
  偶 \((\mathrm{R}(\mathrm{G}),\mathcal{B})\) 是一个自然的含幺基代数（参见 III, p.~130 的习题 6）。
\item
  假设 G 有限。代数 R(G) 这时是有限秩的；计算 \(p_f(\pi)\) （对每个 \(\pi \in \widehat{G}\) ），并辨识与上述引文问题 j) 中定义的元素 r₀ 相对应的 G 的表示。
\end{enumerate}

\begin{enumerate}
\def\labelenumi{\arabic{enumi})}
\setcounter{enumi}{1}
\tightlist
\item
  设 \(n \geqslant 1\) ，\(p \in [1, +\infty]\) 。我们给 \(\mathbf{R}^n\) 赋予范数
\end{enumerate}

\[
\| x \| = \left(\sum_ {i = 1} ^ {n} | x _ {i} | ^ {p}\right) ^ {1 / p} \text {   if   } p \neq + \infty ,
\]

而 \(\|x\| = \sup_{i}|x_{i}|\) （若 \(p = +\infty\) ）。设 G 为赋有该范数的度量空间 \(R^{n}\) 的等距群。

\begin{enumerate}
\def\labelenumi{\alph{enumi})}
\item
  设 \(\sigma\) 为 \(\{1,\ldots,n\}\) 的一个置换，\(\varepsilon\in\{-1,1\}^{n}\) 。设 \((e_{i})\) 为 \(\mathbf{R}^{n}\) 的典范基。\(\mathbf{R}^{n}\) 中使 \(e_{i}\mapsto\varepsilon_{i}e_{\sigma(i)}\) 的唯一自同态属于 G。这种形式的等距所成的集合是 G 的一个子群 W；它是 \(\mathfrak{S}_{n}\) 被 \((\mathbf{Z}/2\mathbf{Z})^{n}\) 的一个扩张。
\item
  群 G 是紧的。
\item
  群 G 在 \(R^{n}\) 上的作用是不可约的；群 G 含于二次型 \(\sum x_{i}^{2}\) 的正交群。
\item
  若 \(p \neq 2\) ，则 G = W。
\end{enumerate}

\begin{enumerate}
\def\labelenumi{\arabic{enumi})}
\setcounter{enumi}{2}
\tightlist
\item
  我们回顾，称群 \(\Gamma\) 是剩余有限的，如果 \(\Gamma\) 的有限指标子群的交约化为 \(\{e\}\) （A, I, p.~129, 习题 5, b)）。
\end{enumerate}

对群 \(\Gamma\) ，我们用 \(\widehat{\Gamma}\) 表示射影极限 \(\widehat{\Gamma} = \varprojlim \Gamma / H\) ，其中 \(H\) 取遍 \(\Gamma\) 的按包含排序的有限指标正规子群所成的集合。我们称 \(\widehat{\Gamma}\) 是 \(\Gamma\) 的射有限完备化。

\begin{enumerate}
\def\labelenumi{\alph{enumi})}
\item
  群 \(\Gamma\) 是剩余有限的，当且仅当对每个异于单位元的 \(g \in \Gamma\) ，存在有限群 F 与同态 \(\varphi: \Gamma \to F\) ，使得 \(\varphi(g) \neq e\) 。
\item
  证明 \(\Gamma\) 是剩余有限的，当且仅当从 \(\Gamma\) 到 \(\widehat{\Gamma}\) 的典范同态是单射。
\item
  设 K 为域。\(\mathbf{GL}(n,\mathrm{K})\) 的每个有限生成子群（参见 A, I, p.~145, 习题 16）都是剩余有限的（“马尔采夫定理”）。（证明存在整数 \(m\geqslant 0\) 、一个等于 \(\mathbf{Z}\) 或某个有限域的环 A，以及该子群到 \(\mathbf{GL}(m,\mathrm{R})\) 的一个单射同态，其中 R 形如 \(\mathrm{A}[\mathrm{X}_1,\dots ,\mathrm{X}_r,1 / f]\) 。）
\item
  设 \(\Gamma\) 为有限生成的离散群。假设存在从 \(\Gamma\) 到某个紧拓扑群的单射同态。对每个异于单位元的 \(\gamma \in \Gamma\) ，存在 \(\Gamma\) 的一个有限维不可约酉表示，使得 \(\gamma\) 在 \(\pi\) 下的像不是恒等映射。
\item
  有限生成群是剩余有限的，当且仅当它同构于某个紧群的一个子群。
\end{enumerate}

¶ 4) 设 G 为群，A 与 B 为 G 的子群。设 \(\varphi\colon A\to B\) 是一个同构。

设 \(r'\) 为自由群 F(G) 中的一个族，使得 G 同构于群 F(G, \(r'\) ) （A, I, p.~87, n°6）。我们称 G 依 \(\varphi\) 而得的 HNN 扩张为群 \(G^{\varphi} = F(X, r)\) ，其中 X 是集合 G、A、B 与一个元素 t 之和，而 r 是 \(r'\) 与关系子 \(t^{-1}at = \varphi(a)\) （对所有 \(a \in A\) ）的和族。

我们用 \(j\colon G\to G^{\varphi}\) 表示唯一的同态，使得 \(j(g)\) 是 \(g\) 在 \(G^{\varphi}\) 中所代表的元素。

\begin{enumerate}
\def\labelenumi{\alph{enumi})}
\tightlist
\item
  设 \(k \geqslant 1\) 为整数，\((g_{0}, \varepsilon_{1}, g_{1}, \ldots, \varepsilon_{k}, g_{k})\) 为一个有限族，使得对每个 i 有 \(g_{i} \in G\) 且 \(\varepsilon_{i} \in \{-1, 1\}\) 。假设不存在满足 \(1 \leqslant j \leqslant k - 1\) 的整数 j 使得
\end{enumerate}

\[
(\varepsilon_ {j}, \varepsilon_ {j + 1}) = (- 1, 1), \quad g _ {i} \in A
\]

或

\[
(\varepsilon_ {j}, \varepsilon_ {j + 1}) = (1, - 1), \quad g _ {i} \in \mathrm{B}.
\]

则元素

\[
j (g _ {0}) t ^ {\varepsilon_ {1}} j (g _ {1}) \dots t ^ {\varepsilon_ {k}} j (g _ {k})
\]

（G \(^{φ}\) 的这个元素）不是单位元（“布里顿引理”。）

\begin{enumerate}
\def\labelenumi{\alph{enumi})}
\setcounter{enumi}{1}
\item
  同态 j 是单射。
\item
  设 \(a\) 与 \(b\) 为不同的符号。设 G 为群 \(\mathrm{F}(\{a,b\}, a^{-1}b^{2}ab^{-3})\) 。记 \(b_{1} = a^{-1}ba\) 。设 H 为 G 的一个有限指标正规子群，\(\pi: \mathrm{G} \to \mathrm{G}/\mathrm{H}\) 为典范投影。有 \([b_{1}, b] \in \mathrm{Ker}(\pi)\) 。
\item
  群 G 不容许忠实的有限维表示。（证明 G 不是剩余有限的。）
\end{enumerate}

\begin{enumerate}
\def\labelenumi{\arabic{enumi})}
\setcounter{enumi}{4}
\tightlist
\item
  设 G 为紧群，\(\varrho\) 为 G 的一个忠实的有限维酉表示。设 \(\pi \in \widehat{G}\) 。
\end{enumerate}

\begin{enumerate}
\def\labelenumi{\alph{enumi})}
\tightlist
\item
  存在整数 a 与 b，使得 \(\pi\) 同构于 \(\varrho^{\otimes a} \otimes \overline{\varrho}^{\otimes b}\) 的一个子表示。（例如可以利用魏尔斯特拉斯–斯通定理。）
\end{enumerate}

我们用 \(w(\pi)\) 表示 \(a + b\) 当 \((a, b)\) 取遍满足这一性质的整数组所成的集合时的最小值。

\begin{enumerate}
\def\labelenumi{\alph{enumi})}
\setcounter{enumi}{1}
\tightlist
\item
  设 \(\Lambda\) 为 \(G^{0}\) 的权格，q 为 \(\Lambda\otimes \mathbf{R}\) 上的一个范数。存在实数 C 与 D，使得对每个表示 \(\pi\in\widehat{G}\) （其在 \(G^{0}\) 上的限制具有最高权 \((\lambda_{1},\ldots,\lambda_{k})\) ），我们有
\end{enumerate}

\[
w (\pi) \leqslant \mathrm{C} \sup (q (\lambda_ {1}), \dots , q (\lambda_ {k})) + \mathrm{D}.
\]

（先考虑连通半单群的情形，并对 G 的适当选取的有限表示族应用 a)。）

\begin{enumerate}
\def\labelenumi{\arabic{enumi})}
\setcounter{enumi}{5}
\tightlist
\item
  设 G 为紧群，\(\varrho\) 为 G 的一个不可约表示。存在中心函数 \(f \in \Theta(G^{\sharp})\) ，使得
\end{enumerate}

\[
f = \sum_ {\pi \in \widehat {\mathrm{G}}} c (\pi) \chi_ {\pi}
\]

其中整数 \(c(\pi)\) 满足

\begin{enumerate}
\def\labelenumi{(\roman{enumi})}
\item
  有 \(c(\pi) \geqslant 0\) （对所有 \(\pi \in \widehat{G}\) ）；
\item
  \(c(1) \leqslant c(\varrho)\) ，且除 \(\varrho\) 是一维的并满足 \(\varrho^{2} = 1\) 的可能情形外，该不等式是严格的；
\item
  \(\mathcal{R}(f) \geqslant 0.\)
\end{enumerate}

¶ 7) 对每个整数 \(n \geqslant 1\) ，我们用 \(\mu_n\) 表示 \(\mathbf{C}\) 上群 \(\mathbf{U}(\mathbf{C}^n) \subset \mathrm{GL}(\mathbf{C}^n)\) 的规范化哈尔测度在迹映射 \(g \mapsto \mathrm{Tr}(g)\) 下的像；它是总质量为 1 的紧支集测度。

\begin{enumerate}
\def\labelenumi{\alph{enumi})}
\tightlist
\item
  设 \(a\) 与 \(b\) 为自然数，满足 \(a + b \leqslant n\) 。我们有
\end{enumerate}

\[
\int_ {\mathbf {C}} z ^ {a} \overline {{z}} ^ {b} d \mu_ {n} (z) = \left\{ \begin{array}{l l} 0 & \text{若 } a \neq b \\ a! & \text{若 } a = b. \end{array} \right.
\]

\begin{enumerate}
\def\labelenumi{\alph{enumi})}
\setcounter{enumi}{1}
\tightlist
\item
  测度序列 \((\mu_{n})_{n\geqslant1}\) 狭收敛于实向量空间 C 上协方差为 \(\mathrm{Q}(z)=|z|^{2}\) 的高斯测度。
\end{enumerate}

\begin{enumerate}
\def\labelenumi{\arabic{enumi})}
\setcounter{enumi}{7}
\item
  设 G 为紧群，\(\varrho_{1}\) 与 \(\varrho_{2}\) 为 G 的不可约酉表示。设 \(n \geqslant 2\) 为整数。若 \(S^{n}\varrho_{1}\) 同构于 \(S^{n}\varrho_{2}\) ，则存在 G 的一个有限指标子群 H，使得 \(\varrho_{1}\) 在 H 上的限制同构于 \(\varrho_{2}\) 在 H 上的限制。
\item
  设 \(k\) 为含 5 个元素的有限域，\(\mathrm{G} = \mathbf{S}\mathbf{L}(2,k)\) 。
\end{enumerate}

\begin{enumerate}
\def\labelenumi{\alph{enumi})}
\tightlist
\item
  群 G 的共轭类由 e、-e 以及七个元素
\end{enumerate}

\[
g _ {3} = \left( \begin{array}{c c} 1 & 1 \\ 0 & 1 \end{array} \right), \qquad g _ {4} = \left( \begin{array}{c c} 1 & 2 \\ 0 & 1 \end{array} \right), \qquad g _ {5} = \left( \begin{array}{c c} - 1 & 1 \\ 0 & - 1 \end{array} \right), \qquad g _ {6} = \left( \begin{array}{c c} - 1 & 2 \\ 0 & - 1 \end{array} \right)
\]

\[
g _ {7} = \left( \begin{array}{c c} 2 & 0 \\ 0 & - 2 \end{array} \right), \qquad g _ {8} = \left( \begin{array}{c c} 0 & 1 \\ - 1 & 1 \end{array} \right), \qquad g _ {9} = \left( \begin{array}{c c} 0 & 1 \\ - 1 & - 1 \end{array} \right).
\]

代表。确定每个共轭类的基数。

\begin{enumerate}
\def\labelenumi{\alph{enumi})}
\setcounter{enumi}{1}
\tightlist
\item
  设 \(\alpha = \frac{1}{2}(1 + \sqrt{5})\) 。存在 G 的一个二维忠实不可约酉表示 \(\pi\) ，其特征标 \(\chi\) 满足 \(\chi(-e) = -2\) 以及
\end{enumerate}

\[
\begin{array}{c} \chi (g _ {4}) = \alpha , \quad \chi (g _ {6}) = - \alpha , \quad \chi (g _ {3}) = - \alpha - 1, \quad \chi (g _ {5}) = \alpha + 1 \\ \chi (g _ {7}) = 0, \quad \chi (g _ {8}) = 1, \quad \chi (g _ {9}) = - 1. \end{array}
\]

（可以考虑 \(\mathbf{GL}(2,k)\) 的四维不可约表示在 G 上的限制，参见 A, VIII, p.~422, 习题 27。）

\begin{enumerate}
\def\labelenumi{\alph{enumi})}
\setcounter{enumi}{2}
\tightlist
\item
  有 \(\mathrm{M}_4(\pi) = 2\) 。
\end{enumerate}

\begin{enumerate}
\def\labelenumi{\arabic{enumi})}
\setcounter{enumi}{9}
\tightlist
\item
  \begin{enumerate}
  \def\labelenumii{\alph{enumii})}
  \tightlist
  \item
    存在一个正交型的不可约表示 \(\pi\) （交错群 \(\mathfrak{A}_5\) 的），使得 \(\mathrm{M}_4(\pi) = 3\) 。
  \end{enumerate}
\end{enumerate}

\begin{enumerate}
\def\labelenumi{\alph{enumi})}
\setcounter{enumi}{1}
\tightlist
\item
  设 R 为 \(E_{6}\) 型（相应地 \(E_{7}\) 型、\(E_{8}\) 型）的根系，W 为 R 的外尔群。设 \(\pi\) 为 W 的几何表示（LIE, V, § 4, 第 3 目）。它是不可约的、正交型的，且满足 \(M_{4}(\pi)=3\) 。
\end{enumerate}

（可以利用诸如 GAP 或 MAGMA 的软件来确定 W 的共轭类及其在 \(\pi\) 下的像的列表。）

\begin{enumerate}
\def\labelenumi{\arabic{enumi})}
\setcounter{enumi}{10}
\item
  设 G 为紧群，\(\pi\) 为 G 在希尔伯特空间 E 中的有限维酉表示。若 \(\mathrm{M}_{4}(\pi) \leqslant 5\) ，则表示 \(\pi\) 是不可约的。
\item
  设 G 为连通紧半单李群，其复化李代数同构于单李代数 \(e_{6}\) 。
\end{enumerate}

\begin{enumerate}
\def\labelenumi{\alph{enumi})}
\item
  存在 G 的一个 27 维忠实不可约表示 \(\varrho\) ；这样的表示是复型的。
\item
  表示 \(\varrho\) 的四阶矩等于 3。
\end{enumerate}

\begin{enumerate}
\def\labelenumi{\arabic{enumi})}
\setcounter{enumi}{12}
\tightlist
\item
  设 G 为连通紧半单李群，g 为其李代数。假设存在 G 的一个 27 维忠实不可约表示 \(\varrho\) 。我们假设 \(\varrho\) 的四阶矩等于 3，且表示 \(\varrho\) 是复型的。
\end{enumerate}

\begin{enumerate}
\def\labelenumi{\alph{enumi})}
\item
  李代数 g 是单的。
\item
  李代数 g 同构于 \(\mathfrak{e}_{6}\) 。
\end{enumerate}

\begin{enumerate}
\def\labelenumi{\arabic{enumi})}
\setcounter{enumi}{13}
\item
  设 G 为连通紧半单李群。假设存在 G 的一个 27 维忠实不可约表示 \(\varrho\) 。G 的复化李代数同构于 \(\mathfrak{e}_6\) ，当且仅当表示 \(\operatorname{End}(\varrho)\) 中存在一个 78 维的不可约子表示。
\item
  设 \(q\) 为奇素数的幂，F 为含 \(q\) 个元素的有限域。考虑由 F 上系数的 3 阶上三角幺幂矩阵组成的海森伯群 H \(_3\) (F)。对 \((a, b, c) \in \mathrm{F}^3\) ，我们记
\end{enumerate}

\[
\sigma (a, b, c) = \left( \begin{array}{c c c} 1 & a & c \\ 0 & 1 & b \\ 0 & 0 & 1 \end{array} \right) \in \mathrm{H} _ {3} (\mathrm{F}).
\]

因此映射 \(\sigma\) 是双射。

\begin{enumerate}
\def\labelenumi{\alph{enumi})}
\item
  \(\mathrm{H}_{3}(\mathrm{F})\) 的中心 C 通过把 \(\sigma(a,b,c)\) 映到 c 的群同态与 F 等同。
\item
  存在 \(q^{2}\) 个 \(\mathrm{H}_{3}(\mathrm{F})\) 的一维酉特征标；这些特征标在 C 上是平凡的。
\item
  设 X 为 F 中由元素 \(\sigma(0,b,c)\) （\((b,c)\in\mathrm{F}^{2}\) ）组成的指标为 q 的子群。子群 X 与 \(F^{2}\) 等同。设 \(\psi_{1}\) 与 \(\psi_{2}\) 为 F 的酉特征标，令 \(\psi(\sigma(0,b,c))=\psi_{1}(b)\psi_{2}(c)\) 。酉表示 \(\varrho\) （\(\mathrm{H}_{3}(\mathrm{F})\) 的、由 X 的特征标 \(\psi\) 诱导）是不可约的，当且仅当 \(\psi_{2}\) 非平凡；它在同构意义下只依赖于 \(\psi_{2}\) ，其中心特征标等于 \(\psi_{2}\) 。\(\varrho\) 的特征标的支集等于 C。
\item
  在前面两个问题中构造的不可约表示就是 \(\mathrm{H}_{3}(\mathrm{F})\) 的全部不可约酉表示。
\end{enumerate}

\begin{enumerate}
\def\labelenumi{\alph{enumi})}
\setcounter{enumi}{4}
\tightlist
\item
  设 \(\varrho\) 为 \(\mathrm{H}_{3}(\mathrm{F})\) 的 q 维不可约酉表示之一。表示 \(\varrho\) 是忠实的，且满足 \(\mathrm{M}_{4}(\varrho)=q^{2}\) 。
\end{enumerate}

\begin{enumerate}
\def\labelenumi{\arabic{enumi})}
\setcounter{enumi}{15}
\tightlist
\item
  设 \(r \geqslant 1\) 为整数，G 为 \(\mathbf{U}(\mathbf{C}^{r})\) 的连通紧子群，它在 \(C^{r}\) 上不可约地作用。我们用 D 表示 \(\mathbf{U}(\mathbf{C}^{r})\) 中由对角矩阵组成的子群。
\end{enumerate}

\begin{enumerate}
\def\labelenumi{\alph{enumi})}
\tightlist
\item
  设 A 为 D 的一个含于 G 在 \(\mathbf{U}(\mathbf{C}^{r})\) 中的正规化子内的子群。设 S 为 \(\widehat{A}\) 中由对角系数 \(a = (a_{i,j}) \mapsto a_{i,i}\) （\(1 \leqslant i \leqslant r\) ）组成的子集。假设 S 含有 r 个元素，且方程
\end{enumerate}

\[
\chi_ {1} \chi_ {2} ^ {- 1} = \chi_ {3} \chi_ {4} ^ {- 1}
\]

（其中 \(\chi_{i} \in S\) ）只有形如 \((\chi_{1}, \chi_{1}, \chi_{3}, \chi_{3})\) 与 \((\chi_{1}, \chi_{2}, \chi_{1}, \chi_{2})\) 的解（这时称 \(S\) 是一个 Sidon 集）。

若 G 是半单的，证明 \(\mathbf{G} = \mathbf{S}\mathbf{U}(\mathbf{C}^{r})\) 。（考虑 A 在 \(\operatorname{End}(\mathbf{C}^{r})\) 上的表示，并由假设推出 G 的复化李代数在 D 的共轭下稳定。）

\begin{enumerate}
\def\labelenumi{\alph{enumi})}
\setcounter{enumi}{1}
\tightlist
\item
  群 G 含有 \(\mathbf{SU}(\mathbf{C}^{r})\) ，当且仅当存在元素 \(g \in G\) ，使得 g 的特征值互不相同并在 \(C^{*}\) 中构成一个 Sidon 集。
\end{enumerate}

\begin{enumerate}
\def\labelenumi{\arabic{enumi})}
\setcounter{enumi}{16}
\tightlist
\item
  设 G 为紧群，H 为 G 的闭子群。设 \(\varrho\) 为 H 在有限维希尔伯特空间 E 中的酉表示。
\end{enumerate}

\begin{enumerate}
\def\labelenumi{\alph{enumi})}
\tightlist
\item
  设 \(\pi \in \widehat{\mathbf{G}}\) 。空间 \(\mathrm{Hom}_{\mathrm{G}}(\pi, \mathrm{Ind}_{\mathrm{H}}^{\mathrm{G}}(\varrho))\) 是有限维的。

\end{enumerate}

\begin{enumerate}
\def\labelenumi{\alph{enumi})}
\setcounter{enumi}{1}
\item
  对每个 \(u \in \mathrm{Hom}_{\mathrm{G}}(\pi, \mathrm{Ind}_{\mathrm{H}}^{\mathrm{G}}(\varrho))\) ，\(u\) 的像含于 \(\mathcal{C}(\mathrm{G};\mathrm{E})\) 在 \(\mathrm{L}^2 (\mathrm{G};\mathrm{E})\) 中的像。
\item
  存在线性映射 \(\Phi\) ，它从空间 \(\mathrm{Hom}_{\mathrm{G}}(\pi, \mathrm{Ind}_{\mathrm{H}}^{\mathrm{G}}(\varrho))\) 到空间 \(\mathrm{Hom}_{\mathrm{H}}(\mathrm{Res}_{\mathrm{H}}^{\mathrm{G}}(\pi), \varrho)\) ，使得对每个 \(u \in \mathrm{Hom}_{\mathrm{G}}(\pi, \mathrm{Ind}_{\mathrm{H}}^{\mathrm{G}}(\varrho))\) 与每个 \(x \in \mathrm{E}_{\pi}\) ，\(\Phi(u)(x)\) 是连续函数在 \(e\) 处的值，该连续函数表示 \(u(x)\) 。
\item
  线性映射 \(\Phi\) 是同构。
\item
  假设 \(\varrho\) 不可约。\(\pi\) 在 \(\mathrm{Ind}_{\mathrm{H}}^{\mathrm{G}}(\varrho)\) 中的重数等于 \(\varrho\) 在 \(\mathrm{Res}_{\mathrm{H}}^{\mathrm{G}}(\pi)\) 中的重数（弗罗贝尼乌斯互反公式，参见 A, VIII, p. 392）。
\end{enumerate}

\begin{enumerate}
\def\labelenumi{\arabic{enumi})}
\setcounter{enumi}{17}
\item
  假设 G 是紧的。设 \(\varrho\) 与 \(\varrho'\) 为 G 的酉表示。表示 \(\varrho\) 弱包含于 \(\varrho'\) （参见 V, p. 497 的习题 11），当且仅当对 G 的每个不可约酉表示 \(\pi\) ，条件 \(m_{\pi}(\varrho) \geqslant 1\) 蕴含 \(m_{\pi}(\varrho') \geqslant 1\) 。（化归到 \(\varrho'\) 不可约的情形；对每个规范化的对角矩阵系数 \(f\) （属于 \(\varrho\) ），应用 V, p.~498 的习题 13 证明存在 \(\varrho'\) 的一个与 \(f\) 不正交的对角矩阵系数，再用正交关系得出结论。）
\item
  设 G 为赋有规范化哈尔测度 \(\mu\) 的紧群，\(f \in \mathcal{C}(\mathrm{G})\) 为 G 上的连续函数。存在不可约表示 \(\pi \in \widehat{G}\) ，使得 \(f = \dim(\pi)^{-1}\chi_{\pi}\) ，当且仅当
\end{enumerate}

\[
\int_ {\mathrm{G}} f (x g x ^ {- 1} h) d \mu (x) = f (g) f (h)
\]

对所有 \((g,h)\in \mathrm{G}^2\) 成立

\begin{enumerate}
\def\labelenumi{\arabic{enumi})}
\setcounter{enumi}{19}
\tightlist
\item
  设 H 为赋有左哈尔测度 \(\nu\) 的局部紧拓扑群。设 G 为 H 的紧子群。设 \(\varrho\) 为 H 在希尔伯特空间 E 中的不可约酉表示。
\end{enumerate}

我们假设 \(\pi\) 在 \(\varrho\) 到 G 的限制中的重数对每个不可约表示 \(\pi \in \widehat{G}\) 都是有限的。我们把 G 上的函数与 H 上紧支集的测度等同起来。

\begin{enumerate}
\def\labelenumi{\alph{enumi})}
\item
  对每个 \(\pi \in \widehat{\mathrm{G}}\) ，E 的自同态 \(\varrho(f * \chi_{\pi})\) 是有限秩的。
\item
  对每个 \(f \in \mathrm{L}^{1}(\mathrm{H})\) ，自同态 \(\varrho(f)\) 是紧的。
\end{enumerate}

\begin{enumerate}
\def\labelenumi{\arabic{enumi})}
\setcounter{enumi}{20}
\tightlist
\item
  设 \(n \geqslant 1\) 为整数。
\end{enumerate}

\begin{enumerate}
\def\labelenumi{\alph{enumi})}
\tightlist
\item
  对每个整数 \(r > n^{2}\) 与每个族 \((a_{1}, \ldots, a_{r})\) （\(\mathcal{L}(\mathbf{C}^{n})\) 的元素），我们有
\end{enumerate}

\[
\sum_ {\sigma \in \mathfrak {S} _ {r}} \varepsilon (\sigma) a _ {\sigma (1)} \dots a _ {\sigma (r)} = 0.\tag{4}
\]

\begin{enumerate}
\def\labelenumi{\alph{enumi})}
\setcounter{enumi}{1}
\tightlist
\item
  我们用 \(r(n)\) 表示最小的整数 \(r \geqslant 1\) ，使公式 (4) 对每个族 \((a_i)_{1 \leqslant i \leqslant r}\) 成立。我们有 \(r(n) \geqslant r(n-1) + 2\) 。（考虑一族 \((a_1, \ldots, a_{r(n-1)-1})\) （属于 \(\mathcal{L}(\mathbf{C}^{n-1})\) ），使得记作 b 的上述表达式非零；把 \(\mathcal{L}(\mathbf{C}^n)\) 的元素解释为矩阵，选取 i 与 j 使 \(b_{i,j}\) 非零；把诸 \(a_i\) 与 b 视为 \(\mathcal{L}(\mathbf{C}^n)\) 的元素（通过添加一行一列零），然后定义矩阵 \(a_{r(n-1)}\) 与 \(a_{r(n-1)+1}\) ，使得 \((a_1, \ldots, a_{r(n-1)+1})\) 的表达式 (4) 非零。）
\end{enumerate}

\begin{enumerate}
\def\labelenumi{\arabic{enumi})}
\setcounter{enumi}{21}
\tightlist
\item
  设 \(r \geqslant 1\) 为整数，\(\mathrm{H} \subset \mathbf{GL}(r, \mathbf{C})\) 为闭子群。
\end{enumerate}

\begin{enumerate}
\def\labelenumi{\alph{enumi})}
\item
  H 上形如 \(g \mapsto \langle \varrho(g)x, \lambda \rangle\) 的函数所成的集合（其中 \(\varrho\) 是 H 在有限维复向量空间 E 中的连续线性表示，\(x \in E\) ，\(\lambda \in E'\) ）是 \(\mathcal{C}(\mathrm{H})\) 的稠密对合子代数。
\item
  对任意非零的函数 \(f \in \mathcal{K}(\mathrm{H})\) ，存在 H 在有限维复向量空间中的一个连续线性表示 \(\varrho\) ，使得 \(\varrho(f)\) 非零。
\item
  假设 H 是连通半单实李群。对每个非零的 \(f \in \mathcal{K}(H)\) ，存在 H 在有限维复向量空间中的不可约连续线性表示 \(\varrho\) ，使得 \(\varrho(f)\) 非零。
\end{enumerate}

\begin{enumerate}
\def\labelenumi{\arabic{enumi})}
\setcounter{enumi}{22}
\tightlist
\item
  我们沿用上一习题的假设与记号，并假设 H 是连通半单实李群。
\end{enumerate}

设 G 为 H 的紧子群。假设存在 H 的连通可解闭子群 N，使得 H = GN。（对每个半单实李群，当 G 是 H 的极大紧子群时这一条件满足；例如当 G = SL(n, R) 时，参见 INT, VII, p.~91, § 3, n° 3, 命题 7。）

设 \(\pi \in \widehat{\mathbf{G}}\) ，\(e_{\pi} = \dim (\pi)\overline{\chi}_{\pi}\) 。设 \(\varrho\) 为上一习题最后一问给出的 H 的不可约表示；我们用 E 表示 \(\varrho\) 的空间。

\begin{enumerate}
\def\labelenumi{\alph{enumi})}
\item
  \(\varrho\) 到 G 的限制存在循环向量。（利用李定理，LIE, III, p.~241, 推论，找出非零的 \(x \in E\) ，使 C \(x\) 在 N 的作用下不变。）\(\pi\) 在 \(\varrho\) 到 G 的限制中的重数至多等于 \(\dim(\pi)\) 。
\item
  设 \(r\) 为习题 22 中定义的整数 \(r(\dim (\pi)^2)\) ；恒等式
\end{enumerate}

\[
\sum_ {\sigma \in \mathfrak {S} _ {r}} \varepsilon (\sigma) a _ {\sigma (1)} \dots a _ {\sigma (r)} = 0
\]

对每个元素族 \(a_{i}=e_{\pi}*g_{i}*e_{\pi}\) （属于 H 上紧支集测度的代数 \(\mathcal{M}_{c}(\mathrm{H})\) ，其中 \(g_{i}\in\mathcal{K}(\mathrm{H})\) （对 \(1\leqslant i\leqslant r\) ））成立。

\begin{enumerate}
\def\labelenumi{\alph{enumi})}
\setcounter{enumi}{2}
\tightlist
\item
  设 \(\varpi\) 为 H 在希尔伯特空间 F 中的不可约酉表示。设 \(F_{\pi}\) 为 \(\pi\)-同型分量（在 \(\varpi\) 到 G 的限制中），\(p\) 为 F 的以 \(F_{\pi}\) 为像的正交投影算子。对每个族 \((u_i)_{1 \leqslant i \leqslant r}\) （属于 \(\mathcal{L}(F_{\pi})\) ），我们有
\end{enumerate}

\[
\sum_ {\sigma \in \mathfrak {S} _ {r}} \varepsilon (\sigma) u _ {\sigma (1)} \dots u _ {\sigma (r)} = 0
\]

  （利用 V, p.~486 的习题 1）。

\begin{enumerate}
\def\labelenumi{\alph{enumi})}
\setcounter{enumi}{3}
\tightlist
\item
  \(\pi\) 在 \(\varpi\) 到 G 的限制中的重数至多等于 \(\dim(\pi)\) 。（利用 V, p.~510 的习题 22 的结果。）
\end{enumerate}

\begin{enumerate}
\def\labelenumi{\arabic{enumi})}
\setcounter{enumi}{23}
\tightlist
\item
  称拓扑群 G 是极小概周期的，如果从 G 到紧群的每个连续同态都是平凡的。
\end{enumerate}

\begin{enumerate}
\def\labelenumi{\alph{enumi})}
\item
  拓扑群 G 是极小概周期的，当且仅当 G 不容许非平凡的有限维酉表示。若 G 是交换的，则 G 是极小概周期的，当且仅当 G 的每个连续酉特征标都是平凡的。
\item
  设 G 为拓扑群，\(g \in G\) 满足：对每个整数 \(k \geqslant 1\) ，存在整数 \(m \geqslant 1\) ，使得 \(g^{km}\) 共轭于 g。则有 \(\varrho(g) = 1_{\mathrm{E}}\) 对 G 在希尔伯特空间 E 中的每个有限维酉表示 \(\varrho\) 成立。
\item
  设 k 为 C 的子域。赋有离散拓扑的群 \(\mathbf{SL}(2,k)\) 是极小概周期的。（验证上一问题适用于 \(\begin{pmatrix}1 & z \\ 0 & 1\end{pmatrix}\) （对所有 \(z \in k\) ）。）
\end{enumerate}

\begin{enumerate}
\def\labelenumi{\arabic{enumi})}
\setcounter{enumi}{24}
\tightlist
\item
  设 G 为赋有左哈尔测度 \(\nu\) 的局部紧群，\(K \subset G\) 为赋有规范化哈尔测度 \(\mu\) 的 G 的紧子群。
\end{enumerate}

我们用 \(\mathcal{H}(\mathrm{G},\mathrm{K})\) 表示由函数 \(\varphi \in \mathcal{H}(\mathrm{G})\) （满足 \(\varphi (k_1gk_2) = \varphi (g)\) ，对所有 \(g\in \mathrm{G}\) 与每个 \((k_{1},k_{2})\in \mathrm{K}\times \mathrm{K}\) ）所成的空间。

空间 \(\mathcal{H}(\mathrm{G},\mathrm{K})\) 是代数 \(\mathcal{M}_c(\mathrm{G})\) 的一个含幺子代数；我们称 \((\mathrm{G},\mathrm{K})\) 为盖尔范德对，如果代数 \(\mathcal{H}(\mathrm{G},\mathrm{K})\) 是交换的。\(a)\) 映射 \(u^{\mathrm{K}}\) 把 \(f\) 对应到函数 \(f^{\mathrm{K}}\) （由

\[
f ^ {\mathrm{K}} (g) = \int_ {\mathrm{K} \times \mathrm{K}} f (k _ {1} g k _ {2}) d (\mu \otimes \mu) (k _ {1}, k _ {2})
\]

所定义）；它是 \(\mathcal{K}(\mathrm{G})\) 中的一个投影算子，其像为 \(\mathcal{H}(\mathrm{G},\mathrm{K})\) 。

\begin{enumerate}
\def\labelenumi{\alph{enumi})}
\setcounter{enumi}{1}
\tightlist
\item
  设 \(f\in \mathcal{K}(\mathrm{G})\) 。我们有
\end{enumerate}

\[
\int_ {\mathrm{G}} f (g) d \nu (g) = \int_ {\mathrm{G}} f ^ {\mathrm{K}} (g) d \nu (g)
\]

并且 \((f\varphi)^{\mathrm{K}} = f^{\mathrm{K}}\varphi\) 对每个函数 \(\varphi \in \mathcal{H}(\mathrm{G},\mathrm{K})\) 成立。

\begin{enumerate}
\def\labelenumi{\alph{enumi})}
\setcounter{enumi}{2}
\item
  设 \(C \subset G\) 为紧子集。存在函数 \(\varphi \in \mathcal{H}(G, K)\) ，使得 \(\varphi(x) = 1\) 对所有 \(x \in C\) 成立。
\item
  设 (G,K) 为盖尔范德对。群 G 是幺模的。（注意 G 的模属于 \(\mathcal{H}(\mathrm{G},\mathrm{K})\) ，并把 \(\varphi \in \mathcal{H}(\mathrm{G},\mathrm{K})\) 的积分解释为适当选取的卷积在 \(e\) 处的值。）
\item
  假设存在连续对合 \(\theta\colon\mathrm{G}\to\mathrm{G}\) ，使得 \(\mathrm{K}\theta(x)\mathrm{K}=\mathrm{K}x^{-1}\mathrm{K}\) 对所有 \(x\in\mathrm{G}\) 成立。则 (G,K) 是盖尔范德对。（验证映射 \(f\mapsto\nu(f\circ\theta)\) （在 \(\mathcal{K}(\mathrm{G})\) 上）是 G 上的一个左哈尔测度，进而它等于 \(\underline{\nu}\) ；推出 \((f*g)\circ\theta=(f\circ\theta)*(g\circ\theta)\) ，并与公式 \(f*g=\check{g}*f\) 比较。）
\item
  若 G 是紧的，则 \((\mathbf{G} \times \mathbf{G}, \Delta)\) 是盖尔范德对，其中 \(\Delta\) 是 \(G \times G\) 的对角子群，即同态 \(x \mapsto (x, x)\) （从 G 到 \(G \times G\) ）的像。
\end{enumerate}

\begin{enumerate}
\def\labelenumi{\arabic{enumi})}
\setcounter{enumi}{25}
\tightlist
\item
  我们继续沿用上一习题的记号。我们用 \(\varphi_{\mathrm{K}}\) 表示 K 的特征函数。
\end{enumerate}

\begin{enumerate}
\def\labelenumi{\alph{enumi})}
\item
  设 \(\pi\) 为 G 在复希尔伯特空间 E 中的酉表示。\(\mathcal{M}_c(\mathrm{G})\) 在 E 中的表示通过过渡到子空间，定义了代数 \(\mathcal{H}(\mathrm{G},\mathrm{K})\) 在 E 的 K 不变向量所成空间 \(\mathrm{E}^{\mathrm{K}}\) 中的表示。
\item
  若 E 容许循环向量 x，使得 \(x \in E^{K}\) ，则 \(\pi\) 是不可约的。
\item
  假设 \(\pi\) 不可约。若空间 \(E^{K}\) 不约化为 0，则 \(\mathcal{H}(G,K)\) 在 \(E^{K}\) 中的表示是不可约的。
\end{enumerate}

\begin{enumerate}
\def\labelenumi{\alph{enumi})}
\setcounter{enumi}{3}
\tightlist
\item
  设 \(\varphi\) 为 G 上的正定函数，\((\varrho, x)\) 为 \(\varphi\) 的一个循环希尔伯特实现。我们有 \(\varphi \in \mathcal{H}(\mathrm{G}, \mathrm{K})\) ，当且仅当 \(x \in \mathrm{E}^{\mathrm{K}}\) 。
\end{enumerate}

\begin{enumerate}
\def\labelenumi{\arabic{enumi})}
\setcounter{enumi}{26}
\tightlist
\item
  我们继续沿用上一习题的记号。假设 (G, K) 是盖尔范德对。
\end{enumerate}

我们称 \(\varphi\in\mathcal{H}(\mathrm{G},\mathrm{K})\) 是 K-球函数，如果映射

\[
f \mapsto \int_ {\mathrm{G}} f \check {\varphi} d \nu
\]

是代数 \(\mathcal{H}(\mathrm{G},\mathrm{K})\) 的一个非零特征标。

\begin{enumerate}
\def\labelenumi{\alph{enumi})}
\tightlist
\item
  设 \(\varphi \in \mathcal{H}(\mathrm{G},\mathrm{K})\) 。证明下列条件等价：
\end{enumerate}

\begin{enumerate}
\def\labelenumi{(\roman{enumi})}
\item
  函数 \(\varphi\) 是 K-球函数；
\item
  对所有 \(x\) 与 \(y\) （属于 \(\mathrm{G}\) ），我们有
\end{enumerate}

\[
\int_ {\mathrm{K}} \varphi (x k y) d \mu (k) = \varphi (x) \varphi (y);
\]

\begin{enumerate}
\def\labelenumi{(\roman{enumi})}
\setcounter{enumi}{2}
\tightlist
\item
  有 \(\varphi(e)=1\) ，且对每个函数 \(\psi\in\mathcal{H}(\mathrm{G},\mathrm{K})\) ，存在 \(c\in\mathbf{C}\) 使得 \(\psi*\varphi=c\varphi\) 。
\end{enumerate}

\begin{enumerate}
\def\labelenumi{\alph{enumi})}
\setcounter{enumi}{1}
\item
  我们用 \(\mathcal{H}^1 (\mathrm{G},\mathrm{K})\) 表示由 \(f\in \mathrm{L}^1 (\mathrm{G})\) （满足 \(\varrho_{\mathrm{G}}(k_1,k_2)f = f\) ，对所有 \((k_{1},k_{2})\in \mathrm{K}\times \mathrm{K}\) ）所成的空间。它是对合巴拿赫代数 \(\mathrm{L}^1 (\mathrm{G})\) 的一个巴拿赫子代数；它是交换的。
\item
  代数 \(\mathcal{H}(\mathrm{G},\mathrm{K})\) 与 \(\mathcal{H}^1 (\mathrm{G},\mathrm{K})\) 的一个稠密子代数等同。\\
\item
  把 \(\varphi \in \mathcal{H}(\mathrm{G},\mathrm{K})\) 对应到线性型
\end{enumerate}

\[
f \mapsto \int_ {\mathrm{G}} f \check {\varphi} d \nu
\]

的映射（在 \(\mathcal{H}^{1}(\mathrm{G},\mathrm{K})\) 上）定义了从 G 上有界 K-球函数所成的空间到 \(\mathcal{H}^{1}(\mathrm{G},\mathrm{K})\) 的非零特征标所成空间上的一个双射。（注意 \(\mathcal{H}^{1}(\mathrm{G},\mathrm{K})\) 的每个特征标都是连续的且范数 \(\leqslant 1\) ，参见 I, p.~104 的命题 2。）

\begin{enumerate}
\def\labelenumi{\arabic{enumi})}
\setcounter{enumi}{27}
\tightlist
\item
  设 G 为局部紧群，K 为 G 的紧子群。
\end{enumerate}

\begin{enumerate}
\def\labelenumi{\alph{enumi})}
\item
  偶 (G, K) 是盖尔范德对，当且仅当对 G 在复希尔伯特空间 E 中的每个酉表示 \(\pi\) ，\(E^{K}\) 的维数至多为 1。（为证该条件是必要的，证明交换对合巴拿赫代数在希尔伯特空间中的不可约表示是一维的；为证它是充分的，利用 V, p.~454 的定理 4 的推论验证代数 \(\mathcal{H}(G,K)\) 的特征标分离其点。）
\item
  假设 (G, K) 是盖尔范德对。设 \(\varphi \in \mathcal{H}(\mathrm{G},\mathrm{K})\) 。假设 \(\varphi\) 是正定函数且 \(\varphi(e) = 1\) 。证明 \(\varphi\) 是 K-球函数，当且仅当 \(\varphi\) 有一个希尔伯特表示 \((\varrho, x)\) ，使得 \(\varrho\) 不可约。
\item
  假设 \((\mathrm{G},\mathrm{K})\) 是盖尔范德对。G 上正定 K-球函数所成的集合与由不可约表示 \(\pi\in\widehat{G}\) （满足 \(E_{\pi}^{K}\) 非零）所成的集合双射对应。
\end{enumerate}

\begin{enumerate}
\def\labelenumi{\arabic{enumi})}
\setcounter{enumi}{28}
\tightlist
\item
  设 G 为紧群，K 为 G 的闭子群。假设 (G, K) 是盖尔范德对。我们用 \(\nu\) （相应地 \(\mu\) ）表示 G（相应地 K）上的规范化哈尔测度。
\end{enumerate}

设 \(\varphi\) 为 G 上的一个 K-球函数。

\begin{enumerate}
\def\labelenumi{\alph{enumi})}
\item
  我们用 \(E_{\varphi}\) 表示 \(\mathcal{C}(G)\) 中形如 \(\alpha * \varphi\) 的函数所成的子空间，其中 \(\alpha\) 是 G 上有限支集的测度。对每个 \(\psi \in \mathcal{H}(G, K)\) ，自同态 \(v_{\psi}: f \mapsto f * \psi\) （\(\mathcal{C}(G)\) 的）是紧的，且 \(E_{\varphi}\) 含于 \(v_{\psi}\) 的一个特征空间。
\item
  空间 \(E_{\varphi}\) 是有限维的。它与正则表示 \(\gamma_{G}\) 的 G-有限向量所成空间的一个含 \(\varphi\) 的子空间等同。
\item
  \(E_{\varphi}\) 的 K 不变向量所成的空间是一维的，且 G 在 \(E_{\varphi}\) 中的表示是不可约的。
\end{enumerate}

\begin{enumerate}
\def\labelenumi{\alph{enumi})}
\setcounter{enumi}{3}
\tightlist
\item
  设 \((f_i)_{i\in I}\) 为 \(\mathrm{E}_{\varphi}\) 的规范正交基，\(k\colon \mathbf{G}\times \mathbf{G}\to \mathbf{C}\) 为由

\[
k (x, y) = \sum_ {i \in \mathrm{I}} \overline {{f _ {i} (x)}} f _ {i} (y).
\]

定义的函数。有

\[
f (x) = \int_ {\mathrm{G}} k (x, y) f (y) d \nu (y)
\]

对所有 \(f \in \mathrm{E}_{\varphi}\) 与 \(x \in \mathrm{G}\) 成立。

\end{enumerate}

\begin{enumerate}
\def\labelenumi{\alph{enumi})}
\setcounter{enumi}{4}
\item
  有 \(k(x,y)=\dim(\mathrm{E}_{\varphi})\varphi(y^{-1}x)\) 对所有 \((x,y)\in\mathrm{G}\times\mathrm{G}\) 成立。
\item
  函数 \(\varphi\) 是正定型的。
\item
  设 \(\varphi\) 为 G 上的 K-球函数，\((\pi, x)\) 为 \(\varphi\) 的一个希尔伯特实现。表示 \(\pi\) 是不可约的；我们有
\end{enumerate}

\[
\varphi (x) = \int_ {\mathrm{K}} \chi_ {\pi} (x ^ {- 1} k) d \mu (k)
\]

对所有 \(x \in \mathrm{G}\) 成立，且 \(\int_{\mathrm{G}} |\varphi|^2 d\nu = \dim(\pi)^{-1}\) 。

\begin{enumerate}
\def\labelenumi{\arabic{enumi})}
\setcounter{enumi}{29}
\tightlist
\item
  设 \(n \in N\) 。用 G 表示仿射群 \(\mathrm{A}(\mathbf{R}^{n})\) （空间 \(R^{n}\) 的）中由元素 \(g \in \mathrm{A}(\mathbf{R}^{n})\) （其相伴线性映射 \(g_{0}\) 属于 \(\mathbf{SO}(\mathbf{R}^{n})\) ）所成的子群；它是一个实李群（参见 A, II, p. 131 与 LIE, III, p. 102, 例）。我们把 G 的元素记作 \(g = (g_{0}, a)\) ，其中 \(g_{0} \in \mathbf{SO}(\mathbf{R}^{n})\) ，\(a \in R^{n}\) 。
\end{enumerate}

\begin{enumerate}
\def\labelenumi{\alph{enumi})}
\item
  偶 \((\mathbf{G},\mathbf{SO}(\mathbf{R}^n))\) 是盖尔范德对。
\item
  设 \(\varphi \in \mathcal{H}(\mathrm{G},\mathbf{SO}(\mathbf{R}^n))\) 。用 \(f\colon \mathbf{R}^n\to \mathbf{C}\) 表示由 \(f(r) = \varphi (1_{\mathbf{R}^N},g(e_1))\) 定义的函数。映射 \(u\colon \varphi \mapsto f\) 是从 \(\mathcal{H}(\mathrm{G},\mathrm{K})\) 到 \(\mathcal{K}(\mathbf{R}^n)\) 的一个单射复代数态射；其像是 \(\mathcal{K}(\mathbf{R}^n)\) 中由函数 \(f\) （满足 \(f(g_0(x)) = f(x)\) ，对所有 \(g_0\in \mathbf{SO}(\mathbf{R}^n)\) 与每个 \(x\in \mathbf{R}^n\) ）组成的子空间。
\item
  我们用 \(\Delta\) 表示 \(\mathbf{R}^n\) 上的拉普拉斯算子。函数 \(\varphi \in \mathcal{K}(\mathrm{G})\) 是球函数，当且仅当函数 \(f = u(\varphi)\) 是 \(\mathrm{C}^{\infty}\) 类的、满足 \(f(0) = 1\) ，且存在 \(\lambda \in \mathbf{C}\) 使得 \(\Delta f = \lambda f\) 。（利用 \(\Delta(f \circ g_0) = (\Delta f) \circ g_0\) 对所有 \(g_0 \in \mathbf{SO}(\mathbf{R}^n)\) 成立这一事实。）
\end{enumerate}

\begin{enumerate}
\def\labelenumi{\arabic{enumi})}
\setcounter{enumi}{30}
\tightlist
\item
  设 \(n\) 为整数 \(\geqslant 1\) 。我们用 N 表示 \(\mathbf{R}^n\) 上的欧几里得范数，用 \(\lambda_n\) 表示 \(\mathbf{R}^n\) 上的勒贝格测度。我们用 \(\omega_n\) 表示球面 \(\mathbf{S}_n \subset \mathbf{R}^{n+1}\) 上唯一的、在正交群 \(\mathbf{O}(n+1, \mathbf{R})\) 的作用下不变的、总质量为 1 的正测度（参见 INT, VIII, p.~116, 习题 8）。
\end{enumerate}

我们把 \(R^{n}\) 与射影空间 \(\mathbf{P}_{n}(\mathbf{R})\) 中射影超平面 H 的余集等同起来（参见 TG, VI, p.~15, 命题 5），并把 \(R^{n}\) 上的测度解释为 \(\mathbf{P}_{n}(\mathbf{R})\) 上支集含于 H 的余集中的测度。

对每个 \(g \in \mathbf{GL}(n+1, \mathbf{R})\) ，我们用 \(\widetilde{g}\) 表示由 g 定义的 \(\mathbf{P}_{n}(\mathbf{R})\) 的同胚。

\begin{enumerate}
\def\labelenumi{\alph{enumi})}
\tightlist
\item
  正测度 \(c_{n}\) （在 \(\mathbf{R}^n\) 上）由
\end{enumerate}

\[
c _ {n} = \pi^ {- (n + 1) / 2} \Gamma \Big (\frac {n + 1}{2} \Big) (1 + \mathrm{N} ^ {2}) ^ {- (n + 1) / 2} \lambda_ {n}
\]

定义的（“柯西测度”）是有界的且总质量为 1。若 \(n = 1\) ，它就是参数 1 的稳定测度（参见习题 V, p.~496）。对每个射影超平面 \(\mathrm{H}'\) （\(\mathbf{P}_n(\mathbf{R})\) 的），我们有 \(c_n(\mathrm{H}') = 0\) 。

\begin{enumerate}
\def\labelenumi{\alph{enumi})}
\setcounter{enumi}{1}
\item
  从 \(\mathbf{R}^{n+1} - \{0\}\) 到 \(\mathbf{P}_n(\mathbf{R})\) 的典范映射通过过渡到子空间，定义了连续满射 \(\sigma: \mathbf{S}_n - (\mathbf{R}^n \times \{0\}) \to \mathbf{R}^n\) 。我们有 \(\sigma(\omega_n) = c_n\) 。
\item
  设 \(g \in \mathbf{GL}(n+1, \mathbf{R})\) 。存在 \(a \in \mathbf{GL}(n, \mathbf{R})\) 与 \(b \in R^n\) ，使得像测度 \(\widetilde{g}(c_n)\) 等于 \(c_n\) 在仿射映射 \(x \mapsto a(x) + b\) 下的像。
\end{enumerate}

在下面，我们打算证明这一断言的一个逆命题。

\begin{enumerate}
\def\labelenumi{\alph{enumi})}
\setcounter{enumi}{3}
\tightlist
\item
  设 \(\mu\) 为 \(\mathbf{P}_{n}(\mathbf{R})\) 上的有界正测度，使得每个射影超平面都是 \(\mu\) -可忽略的。由 \(g \in \mathbf{GL}(n+1, \mathbf{R})\) （满足 \(g(\mu) = \mu\) 且 \(|\det(g)| = 1\) ）所成的集合是 \(\mathbf{GL}(n+1, \mathbf{R})\) 的一个紧子群。

\end{enumerate}

\begin{enumerate}
\def\labelenumi{\alph{enumi})}
\setcounter{enumi}{4}
\item
  设 \(\mu\) 为 \(\mathbf{P}_n(\mathbf{R})\) 上总质量为 1 的有界正测度，使得对每个 \(g \in \mathbf{GL}(n + 1, \mathbf{R})\) ，存在仿射映射 \(f\) （\(\mathbf{R}^n\) 的）满足 \(g(\mu) = f(\mu)\) 。每个射影超平面都是 \(\mu\) -可忽略的。
\item
  设 \(K_{\mu}\) 为由 \(g \in \mathbf{GL}(n+1, \mathbf{R})\) （满足 \(g(\mu) = \mu\) ）所成的子群。存在 \(h \in \mathbf{GL}(n+1, \mathbf{R})\) ，使得 \(K_{\mu} \subset h \mathbf{O}(n+1, \mathbf{R}) h^{-1}\) 。
\item
  设 \(K = h^{-1}K_{\mu}h\) 。我们有 \(\mathbf{O}(n + 1, \mathbf{R}) = \mathbf{G} \cdot \mathbf{K}\) ，其中 G 是 \(\mathbf{O}(n + 1, \mathbf{R})\) 中由矩阵
\end{enumerate}

\[
\left( \begin{array}{c c} x & 0 \\ 0 & 1 \end{array} \right)
\]

（\(x\in \mathbf{O}(n,\mathbf{R})\) ）组成的子群

\begin{enumerate}
\def\labelenumi{\alph{enumi})}
\setcounter{enumi}{7}
\tightlist
\item
  测度 \(h^{-1}(\mu)\) 与 \(c_{n}\) 重合。（注意群 K 传递地作用于 \(S_{n}\) ，且像测度 \(\sigma(h^{-1}(\mu))\) 是 \(S_{n}\) 上的一个 K 不变测度。）
\end{enumerate}

\begin{enumerate}
\def\labelenumi{\arabic{enumi})}
\setcounter{enumi}{31}
\tightlist
\item
  设 G 为局部紧群。假设存在 G 的一个既紧又开的子群 H。我们用 \(\mu\) 表示 G 上满足 \(\mu(\mathrm{H}) = 1\) 的唯一的左哈尔测度，用 \(\varphi\) 表示 H 的特征函数。
\end{enumerate}

\begin{enumerate}
\def\labelenumi{\alph{enumi})}
\item
  设 \(\pi\) 为 G 在复希尔伯特空间 E 中的一个平方可积表示。证明自同态 \(\pi(\varphi)\) 是希尔伯特–施密特自同态。
\item
  空间 \(E^{H}\) 是有限维的，且满足
\end{enumerate}

\[
\dim (\mathrm{E} ^ {\mathrm{H}}) \leqslant \frac {1}{c (\pi)}
\]

其中 \(c(\pi)\) 是 \(\pi\) 相对于 \(\{e\}\) 的形式次数。（利用 V, p.~487 的习题 6。）

\begin{enumerate}
\def\labelenumi{\alph{enumi})}
\setcounter{enumi}{2}
\tightlist
\item
  离散无限群没有平方可积表示；若 \(n \geqslant 2\) 为整数，则群 \(\mathbf{SL}(n, \mathbf{Z})\) 没有模中心的平方可积表示。
\end{enumerate}

